% Báo cáo kỹ thuật: Phân tích ý kiến theo khía cạnh (ASQP) trên bộ dữ liệu OATS
% Định dạng theo mẫu báo cáo Project III của Trường CNTT&TT, ĐHBK Hà Nội.
% Biên dịch: xelatex main -> bibtex main -> xelatex main -> xelatex main
%            (hoặc: latexmk -xelatex main.tex)
\documentclass[12pt,a4paper]{report}

% ============================================================================
% PACKAGES
% ============================================================================
\usepackage{fontspec}
\usepackage[vietnamese]{babel}
\usepackage{amsmath,amssymb}
\usepackage{graphicx}
\usepackage{booktabs}
\usepackage{longtable}
\usepackage{array}
\usepackage{tabularx}
\usepackage{listings}
\usepackage{fancyvrb}
\usepackage[table]{xcolor}
\usepackage{hyperref}
\usepackage{geometry}
\usepackage{setspace}
\usepackage{caption}
\usepackage{float}
\usepackage{enumitem}
\usepackage{url}
\usepackage{tcolorbox}
\usepackage{tikz}
\usetikzlibrary{arrows.meta,positioning,shapes.geometric,calc,fit,backgrounds}
\usepackage{pgfplots}
\pgfplotsset{compat=1.18}

\geometry{margin=2.5cm}
\onehalfspacing
\graphicspath{{hinh/}}
\setcounter{tocdepth}{2}
\hypersetup{hidelinks}
\setlist{itemsep=1pt,topsep=3pt}

\newcommand{\enquote}[1]{“#1”}
\newcommand{\code}[1]{\texttt{#1}}
\renewcommand{\lstlistingname}{Đoạn mã}

\definecolor{oatsblue}{RGB}{18,92,173}
\definecolor{oatsgreen}{RGB}{46,140,87}
\definecolor{oatsamber}{RGB}{214,140,20}
\definecolor{oatsred}{RGB}{190,50,50}
\definecolor{oatsgray}{RGB}{110,116,128}
\definecolor{lightblue}{RGB}{228,239,252}
\definecolor{lightgreen}{RGB}{224,244,232}
\definecolor{lightamber}{RGB}{255,241,214}
\definecolor{lightred}{RGB}{252,226,226}
\definecolor{lightgray}{RGB}{242,243,245}

\newtcolorbox{ghichu}[1]{colback=lightgray,colframe=black!55,boxrule=0.4pt,arc=0pt,
  fonttitle=\bfseries,coltitle=black,colbacktitle=black!10,title=#1}

\tikzset{
  box/.style={draw=oatsblue,fill=lightblue,rounded corners=3pt,align=center,font=\small\sffamily,
              minimum height=9mm,inner sep=4pt,text width=#1},
  box/.default=30mm,
  gbox/.style={box=#1,draw=oatsgreen,fill=lightgreen},
  gbox/.default=30mm,
  abox/.style={box=#1,draw=oatsamber,fill=lightamber},
  abox/.default=30mm,
  rbox/.style={box=#1,draw=oatsred,fill=lightred},
  rbox/.default=30mm,
  db/.style={cylinder,shape border rotate=90,aspect=0.18,draw=oatsblue,fill=lightblue,
             align=center,font=\small\sffamily,minimum height=13mm,text width=24mm},
  arr/.style={-{Stealth[length=2.6mm]},thick,color=oatsgray},
  lab/.style={font=\scriptsize\sffamily,color=oatsgray,fill=white,inner sep=1pt},
}

\begin{document}

% ============================================================================
% TRANG BÌA
% ============================================================================
\begin{titlepage}
\begin{center}

\textsc{\large \textbf{ĐẠI HỌC BÁCH KHOA HÀ NỘI}}\\[0.5cm]

{\Large TRƯỜNG CÔNG NGHỆ THÔNG TIN VÀ TRUYỀN THÔNG \\[1cm]}

\includegraphics[scale=0.13]{logo-soict-hust-1.png}\\[1cm]
\textsc{\Large \textbf{Báo cáo kỹ thuật}}\\[0.4cm]

\rule{\linewidth}{0.7mm} \\[0.4cm]
{\huge \bfseries Phân tích ý kiến theo khía cạnh\\[0.2cm] trên bộ dữ liệu OATS:\\[0.2cm] từ pipeline sinh quadruple\\[0.2cm] đến dashboard cho doanh nghiệp\\[0.4cm]}
{\large Đề tài: Generative Aspect-Based Sentiment Analysis on OATS\\[0.3cm]}
\rule{\linewidth}{0.7mm} \\[1cm]

{\large \textbf{Nhóm thực hiện:}}\\[0.6cm]
\begin{tabular}{ll}
\large Nguyễn Đức Mạnh & \large 20235525 \\[0.2cm]
\large Hoàng Văn Nhân & \large 20235542 \\[0.2cm]
\large Lê Nguyễn Phước Thành & \large 20235561 \\[0.2cm]
\large Trần Quang Trọng & \large 20235565 \\[0.6cm]
\end{tabular}

\vfill

{\large \textbf{Hà Nội - Việt Nam}\\[0.2cm] \color{blue!90!black}2026}

\end{center}
\end{titlepage}
\setstretch{1.5}

\pagenumbering{roman}

\tableofcontents
\newpage

\begin{abstract}
Một khách sạn nhận hàng nghìn nhận xét và thấy điểm trung bình giảm, nhưng điểm số đó không cho biết khách đang phàn nàn
về điều gì. Bài toán \textit{Aspect Sentiment Quad Prediction} (ASQP) tách mỗi ý kiến trong câu thành bộ bốn thành phần
\emph{(đối tượng, khía cạnh, cụm ý kiến, cảm xúc)}, nhờ đó có thể tổng hợp theo từng khía cạnh thay vì theo một điểm chung.

Báo cáo trình bày toàn bộ hệ thống nhóm đã xây dựng quanh bộ dữ liệu OATS~\cite{oats2024}: một pipeline nghiên cứu
(đọc dữ liệu, làm sạch, chia tập, huấn luyện, đánh giá, phân tích lỗi) cho hai mô hình sinh GAS~\cite{gas2021} và
Paraphrase~\cite{paraphrase2021} trên ba lĩnh vực; một service HTTP bọc checkpoint đã huấn luyện; và một giao diện web
gồm hai phần: phân tích một câu, và dashboard nhận vào một tệp nhiều nhận xét rồi cho biết mỗi khía cạnh có bao nhiêu
phần trăm nhận xét tiêu cực, kèm các cụm từ chê phổ biến.

Lượt chạy seed 42 trên ba lĩnh vực cho F1 khớp hoàn toàn trên tập test lần lượt là 0{,}227 / 0{,}111 / 0{,}347 với GAS và
0{,}210 / 0{,}144 / 0{,}389 với Paraphrase. Báo cáo chỉ ra một đặc điểm của bộ dữ liệu làm các con số này không so sánh
trực tiếp được với bài báo: trong tập test có sẵn, 80{,}8\% dòng của \code{coursera} và 86{,}7\% dòng của \code{hotels}
không có nhãn nào, trong khi tỷ lệ này ở tập train chỉ khoảng 25\%. Phần kiểm thử gồm 150 test Python và 98 test
frontend; luồng web đã được chạy thật đầu-cuối qua HTTP. Những phần chưa làm --- chạy đa seed, chạy lại theo giao thức
chia tập cấp bài đánh giá, và đo độ chính xác riêng cho cặp (khía cạnh, cảm xúc) mà dashboard dựa vào --- được nêu rõ ở
Chương~\ref{chap:ketluan} thay vì được suy đoán.
\end{abstract}

\chapter*{Danh mục thuật ngữ và viết tắt}
\addcontentsline{toc}{chapter}{Danh mục thuật ngữ và viết tắt}
\begin{tabularx}{\linewidth}{@{}lX@{}}
\toprule
\textbf{Viết tắt} & \textbf{Nghĩa} \\ \midrule
ABSA & Aspect-Based Sentiment Analysis: phân tích cảm xúc theo khía cạnh \\
ASQP & Aspect Sentiment Quad Prediction: dự đoán bộ bốn (target, aspect, opinion, sentiment) \\
Quadruple & Bộ bốn thành phần của một ý kiến trong câu \\
Target & Đối tượng được nhắc tới; \code{NULL} nếu chỉ được ngầm hiểu \\
Aspect (category) & Khía cạnh, lấy từ danh mục cố định của từng lĩnh vực, ví dụ \code{rooms cleanliness} \\
Opinion & Cụm từ thể hiện ý kiến; \code{NULL} nếu ngầm hiểu \\
Sentiment & Một trong bốn nhãn: positive, negative, neutral, conflict \\
OATS & Opinion Aspect Target Sentiment, bộ dữ liệu của bài báo~\cite{oats2024} \\
GAS & Generative ABSA, sinh chuỗi theo kiểu trích xuất~\cite{gas2021} \\
Paraphrase & Sinh câu diễn giải tự nhiên cho mỗi quadruple~\cite{paraphrase2021} \\
T5 & Text-to-Text Transfer Transformer, mô hình nền dùng cho cả hai phương pháp~\cite{t5} \\
Exact match & Một dự đoán chỉ đúng khi cả bốn thành phần khớp chính xác \\
Leakage & Rò rỉ: các câu của cùng một bài đánh giá nằm ở nhiều tập khác nhau \\
Fingerprint & Mã băm SHA-256 của tệp dữ liệu, dùng để phát hiện dữ liệu đổi \\
Mock mode & Chế độ giao diện chạy bằng dữ liệu mẫu, không gọi mô hình \\
\bottomrule
\end{tabularx}

\newpage
\pagenumbering{arabic}

% ============================================================================
\chapter{Giới thiệu}

\section{Đặt vấn đề}

Khi một doanh nghiệp nhìn vào điểm đánh giá trung bình, họ biết \emph{có} vấn đề nhưng không biết vấn đề \emph{nằm ở đâu}.
Câu \enquote{Vị trí tuyệt vời, nhưng phòng bẩn và nhân viên thiếu nhiệt tình} vừa khen vừa chê; gộp thành một điểm sao sẽ
mất toàn bộ thông tin hành động được. Phân tích cảm xúc ở mức câu cũng không đủ, vì một câu có thể chứa nhiều ý kiến trái
chiều về nhiều đối tượng khác nhau.

ASQP giải quyết đúng chỗ thiếu đó: mỗi ý kiến được tách thành bộ bốn thành phần. Với câu trên, kết quả mong muốn là hai
bộ: \code{(location, location general, "Great", positive)} và \code{(room, rooms cleanliness, "dirty", negative)}.
Khi đã có các bộ bốn này cho hàng nghìn nhận xét, việc tổng hợp \enquote{khía cạnh nào bị chê nhiều nhất} trở thành phép
đếm, và quan trọng hơn, có thể chỉ ra \emph{vì sao} bị chê thông qua các cụm ý kiến đi kèm.

Bộ dữ liệu OATS~\cite{oats2024} được công bố năm 2024 cho ba lĩnh vực: đánh giá thực phẩm trên Amazon, đánh giá khóa học
trực tuyến Coursera và đánh giá khách sạn TripAdvisor. Điểm khiến OATS khó hơn các bộ ABSA trước đó là nó gán nhãn ở mức
bài đánh giá nhiều câu, giữ cả các câu không mang ý kiến, và cho phép cả đối tượng lẫn cụm ý kiến ở dạng ngầm hiểu.

\section{Mục tiêu}

Nhóm đặt ba mục tiêu, theo thứ tự ưu tiên:

\begin{enumerate}
  \item \textbf{Một pipeline nghiên cứu trung thực.} Chuẩn bị dữ liệu, huấn luyện, đánh giá và phân tích lỗi cho hai mô
        hình sinh trên ba lĩnh vực, trong đó mọi thao tác đều để lại dấu vết kiểm chứng được: nguồn dữ liệu, mã băm, cấu
        hình, danh mục khía cạnh.
  \item \textbf{Kết quả thực nghiệm có điều kiện rõ ràng.} Báo cáo số liệu theo từng cặp (lĩnh vực, mô hình), không gộp
        ba lĩnh vực, không thay số của nhóm bằng số của bài báo, và nói rõ giao thức chia tập nào đã được dùng.
  \item \textbf{Một sản phẩm dùng được.} Giao diện web cho phép phân tích một câu và, quan trọng hơn với doanh nghiệp,
        tải lên một tệp nhiều nhận xét để biết khía cạnh nào đang yếu.
\end{enumerate}

\section{Phạm vi và đóng góp}

\subsection{Phạm vi}

Báo cáo giới hạn ở hai phương pháp sinh là GAS và Paraphrase trên nền T5-base, ba lĩnh vực của OATS và dữ liệu tiếng Anh.
Các phương pháp khác được nhắc tới trong phần khảo sát để đối chiếu, không được cài đặt lại. Hệ thống phục vụ hiện dành
cho chạy nội bộ: không có xác thực người dùng, không có giới hạn tần suất.

\subsection{Đóng góp}

\begin{itemize}
  \item Pipeline đầy đủ với ba giao thức chia tập (\code{clean\_split} theo bài đánh giá, \code{original\_split} căn
        chỉnh với XML, \code{provided\_split} dùng thẳng tệp chia sẵn) và cơ chế kiểm tra dấu vân tay dữ liệu xuyên suốt
        từ chuẩn bị dữ liệu đến đánh giá và so sánh.
  \item Lượt chạy seed 42 hoàn chỉnh: 6 lần huấn luyện, 6 lần đánh giá, 3 lần so sánh cùng lĩnh vực, 6 lần phân tích lỗi.
  \item Phát hiện và ghi nhận một đặc điểm dữ liệu ảnh hưởng trực tiếp tới việc diễn giải kết quả: tỷ lệ dòng không nhãn
        trong tập test chia sẵn của \code{coursera} và \code{hotels} cao bất thường (Mục~\ref{sec:nhanrong}).
  \item Service HTTP và giao diện web, trong đó có dashboard tổng hợp theo khía cạnh ở mức bài đánh giá, kèm quy tắc đếm
        được định nghĩa tường minh và công bố ngay trên giao diện.
  \item Bộ kiểm thử tự động 150 test Python và 98 test frontend, cùng một lần kiểm chứng đầu-cuối chạy thật qua HTTP.
\end{itemize}

\section{Cấu trúc báo cáo}

Chương~\ref{chap:cosoly} trình bày cơ sở lý thuyết và khảo sát các phương pháp liên quan. Chương~\ref{chap:yeucau} phân
tích yêu cầu. Chương~\ref{chap:thietke} mô tả kiến trúc và các quyết định thiết kế. Chương~\ref{chap:trienkhai} đi vào
chi tiết triển khai. Chương~\ref{chap:ketqua} trình bày kết quả thực nghiệm, kiểm thử và các giới hạn đã xác định.
Chương~\ref{chap:ketluan} kết luận và nêu hướng phát triển.

% ============================================================================
\chapter{Cơ sở lý thuyết và khảo sát}
\label{chap:cosoly}

\section{Bài toán ASQP}

Cho một câu $s$, mô hình phải sinh ra tập các bộ bốn
\[
Q(s) = \{(t_i, a_i, o_i, p_i)\}_{i=1}^{n},
\]
trong đó $t_i$ là đối tượng, $a_i$ là khía cạnh lấy từ danh mục cố định của lĩnh vực, $o_i$ là cụm từ thể hiện ý kiến, và
$p_i$ là cảm xúc. Cả $t_i$ và $o_i$ có thể mang giá trị \code{NULL} khi chỉ được ngầm hiểu. Số bộ $n$ có thể bằng 0 với
câu không chứa ý kiến.

Chỉ số đánh giá là precision, recall và F1 theo \emph{khớp hoàn toàn}: một dự đoán chỉ được tính đúng khi cả bốn thành
phần trùng khớp. Đây là tiêu chí nghiêm ngặt --- đoán đúng ba trên bốn thành phần vẫn bị tính sai hoàn toàn --- và là lý
do các con số F1 trong lĩnh vực này thấp hơn nhiều so với bài toán phân loại cảm xúc thông thường.

\section{Bộ dữ liệu OATS}

Bảng~\ref{tab:dulieu} thống kê dữ liệu thô do nhóm đo trực tiếp từ các tệp XML của bộ dữ liệu.

\begin{table}[H]
\centering
\caption{Quy mô dữ liệu OATS, đo bằng chính bộ đọc XML của dự án.}
\label{tab:dulieu}
\begin{tabular}{lrrrr}
\toprule
\textbf{Lĩnh vực} & \textbf{Bài đánh giá} & \textbf{Câu} & \textbf{Câu không nhãn} & \textbf{Số khía cạnh} \\
\midrule
\code{amazon\_ff} & 1.850 & 9.200 & 2.613 (28,4\%) & 13 \\
\code{coursera}   & 1.850 & 8.985 & 2.825 (31,4\%) & 30 \\
\code{hotels}     & 1.847 & 9.793 & 3.182 (32,5\%) & 37 \\
\bottomrule
\end{tabular}
\end{table}

Hai đặc điểm của bộ dữ liệu ảnh hưởng trực tiếp tới thiết kế hệ thống:

\begin{itemize}
  \item \textbf{Phân bố nhãn lệch mạnh.} Khoảng 78--83\% quadruple mang nhãn positive, 15--18\% negative, chỉ 2--4\%
        neutral, và nhãn thứ tư \code{conflict} xuất hiện rất hiếm (5--6 lần trên mỗi lĩnh vực). Hệ quả là mọi con số
        liên quan tới neutral và conflict đều dựa trên rất ít ví dụ.
  \item \textbf{Nhãn gắn ở mức bài đánh giá.} Một bài gồm nhiều câu; nếu chia tập ở mức câu thì các câu của cùng một bài
        sẽ nằm ở cả tập train lẫn test. Dự án vì vậy coi \emph{không rò rỉ ở mức bài đánh giá} là một ràng buộc cần được
        kiểm tra bằng mã, không phải một giả định.
\end{itemize}

\section{Hai phương pháp sinh được cài đặt}

\subsection{Paraphrase}

Phương pháp Paraphrase~\cite{paraphrase2021} biến mỗi quadruple thành một câu tiếng Anh tự nhiên theo khuôn
\enquote{\textit{aspect} is \textit{phrase} because \textit{target} is \textit{opinion}}, trong đó cảm xúc được ánh xạ
thành một tính từ: positive $\rightarrow$ \code{great}, negative $\rightarrow$ \code{bad}, neutral $\rightarrow$ \code{ok}.
Nhiều quadruple trong một câu được nối bằng dấu phân cách \code{[SSEP]}. Đối tượng hoặc ý kiến ngầm hiểu được viết là
\code{it}. Ưu điểm là chuỗi đích gần với phân bố ngôn ngữ mà T5 đã học trong giai đoạn tiền huấn luyện.

OATS có nhãn thứ tư là \code{conflict}, không nằm trong ba tính từ của bài báo gốc. Dự án mở rộng bằng từ \code{mixed} và
ghi rõ đây là lựa chọn của nhóm, không phải thiết lập của bài báo.

\subsection{GAS}

Bài báo GAS~\cite{gas2021} trình bày hai kiểu sinh cho các bài toán ABSA, trong đó có kiểu trích xuất. Bài báo OATS áp
dụng GAS cho ASQP nhưng \emph{không công bố} cú pháp chuỗi đích cụ thể. Dự án vì vậy chọn một biểu diễn trích xuất có
cấu trúc: mảng JSON gồm các bộ bốn chuỗi theo thứ tự cố định.

\begin{ghichu}{Ranh giới cần giữ}
Vì cú pháp đích của GAS-T5 trong bài báo OATS không được công bố, kết quả GAS trong báo cáo này là của \emph{một cách
triển khai theo hướng trích xuất do nhóm chọn}, không phải tái tạo chính xác thí nghiệm của bài báo. Mọi so sánh với
cột GAS-T5 trong bài báo phải kèm theo lưu ý này.
\end{ghichu}

\section{Các hướng tiếp cận khác}

Bài báo OATS còn báo cáo Template-ILO và Template-DLO~\cite{templatedlo2021} --- tăng cường dữ liệu bằng cách hoán vị
thứ tự các thành phần trong chuỗi đích --- và GEN-SCL-NAT~\cite{genscl2023} --- kết hợp học đối lập có giám sát với một
cấu trúc đích giàu thông tin hơn. Cả ba đều cùng họ \enquote{sinh chuỗi bằng T5}, khác nhau chủ yếu ở dạng chuỗi đích và
cách tăng cường dữ liệu. Nhóm không cài đặt lại các phương pháp này; chúng được dùng làm mốc đối chiếu về mặt số liệu.

Ở phía sản phẩm, các công cụ phân tích nhận xét thương mại thường dừng ở mức gán nhãn cảm xúc cho từng chủ đề được phát
hiện tự động. Khoảng trống mà hệ thống này nhắm tới là giữ nguyên \emph{cụm từ ý kiến} do mô hình trích ra, nhờ đó
dashboard trả lời được câu hỏi \enquote{vì sao khía cạnh này bị chê} chứ không chỉ \enquote{khía cạnh này bị chê}.

\section{Công nghệ sử dụng}

\begin{table}[H]
\centering
\caption{Công nghệ chính và vai trò.}
\begin{tabularx}{\linewidth}{@{}lX@{}}
\toprule
\textbf{Thành phần} & \textbf{Công nghệ} \\
\midrule
Pipeline nghiên cứu & Python 3.12, PyTorch, Transformers~\cite{huggingface}, mô hình nền T5-base~\cite{t5} \\
Service dự đoán & FastAPI~\cite{fastapi} + Uvicorn, nạp checkpoint qua chính mã của pipeline \\
Giao diện web & React 19, TypeScript, Vite \\
Kiểm thử & pytest (150 test), Vitest + Testing Library (98 test) \\
\bottomrule
\end{tabularx}
\end{table}

% ============================================================================
\chapter{Phân tích yêu cầu}
\label{chap:yeucau}

\section{Tác nhân}

\begin{itemize}
  \item \textbf{Người nghiên cứu} chạy pipeline, cần kết quả tái lập được và bằng chứng về dữ liệu đã dùng.
  \item \textbf{Người phân tích của doanh nghiệp} có một tệp nhận xét và cần biết khía cạnh nào đang yếu, vì sao, và có
        đủ dữ liệu để kết luận hay không.
  \item \textbf{Người phát triển} nối giao diện với mô hình, cần một giao diện lập trình ổn định và chạy được cả khi
        chưa có checkpoint.
\end{itemize}

\section{Yêu cầu chức năng}

\begin{table}[H]
\centering
\caption{Yêu cầu chức năng chính.}
\begin{tabularx}{\linewidth}{@{}llX@{}}
\toprule
\textbf{Mã} & \textbf{Tác nhân} & \textbf{Mô tả} \\
\midrule
CN1 & Nghiên cứu & Chuẩn bị dữ liệu từ nguồn thô, ghi lại nguồn, mã băm và danh mục khía cạnh \\
CN2 & Nghiên cứu & Huấn luyện một thí nghiệm theo tệp cấu hình, chọn checkpoint tốt nhất theo F1 tập dev \\
CN3 & Nghiên cứu & Đánh giá trên tập test, xuất metrics, dự đoán và các ca lỗi \\
CN4 & Nghiên cứu & So sánh hai mô hình trong cùng một lĩnh vực, có kiểm tra cùng dữ liệu \\
CN5 & Phát triển & Dự đoán một câu qua HTTP \\
CN6 & Phát triển & Dự đoán nhiều câu trong một lần gọi, phục vụ xử lý theo lô \\
CN7 & Phân tích & Phân tích một câu trên giao diện và xem đủ bốn thành phần \\
CN8 & Phân tích & Tải lên tệp nhận xét (CSV/TXT/JSONL/JSON) và xem tỷ lệ cảm xúc theo từng khía cạnh \\
CN9 & Phân tích & Xem các cụm từ chê phổ biến và các nhận xét cụ thể của một khía cạnh \\
CN10 & Phân tích & Xuất bảng tổng hợp ra CSV \\
\bottomrule
\end{tabularx}
\end{table}

\section{Yêu cầu phi chức năng}

\begin{itemize}
  \item \textbf{Trung thực về bằng chứng.} Giao diện không được hiển thị con số chất lượng mô hình khi chưa có bằng
        chứng; chế độ dữ liệu mẫu phải được dán nhãn rõ ràng.
  \item \textbf{Không rò rỉ chi tiết nội bộ.} Lỗi phía máy chủ chỉ hiện thông báo ngắn và mã lỗi; vết lỗi nằm trong log.
  \item \textbf{Chạy được khi chưa có mô hình.} Thiếu checkpoint phải cho ra trạng thái \code{MODEL\_NOT\_READY} có thể
        hiển thị, không phải sự cố.
  \item \textbf{Tiếp cận được.} Nhãn gắn đúng ô nhập, thông báo bằng vùng \code{aria-live}, cảm xúc phân biệt bằng chữ và
        ký hiệu chứ không chỉ bằng màu~\cite{wcag}.
  \item \textbf{Tái lập được.} Mỗi thí nghiệm ghi lại cấu hình, dấu vân tay dữ liệu và commit mã nguồn.
\end{itemize}

% ============================================================================
\chapter{Thiết kế hệ thống}
\label{chap:thietke}

\section{Kiến trúc tổng thể}

Hệ thống gồm ba tầng tách biệt, nối với nhau bằng những ranh giới hẹp và được định nghĩa rõ: tầng nghiên cứu sinh ra
checkpoint, tầng phục vụ bọc checkpoint thành HTTP, tầng giao diện tiêu thụ HTTP (Hình~\ref{fig:kientruc}).

\begin{figure}[H]
\centering
\begin{tikzpicture}[node distance=7mm]
  \node[db] (raw) {Dữ liệu OATS\\(XML, TXT)};
  \node[box=32mm,right=12mm of raw] (prep) {Chuẩn bị dữ liệu\\\scriptsize đọc, làm sạch, chia tập};
  \node[db,right=12mm of prep] (proc) {\code{data/processed}\\\scriptsize jsonl + metadata};
  \node[box=32mm,below=10mm of proc] (train) {Huấn luyện\\\scriptsize T5-base, chọn best};
  \node[db,left=12mm of train] (ckpt) {\code{artifacts/}\\\scriptsize checkpoint};
  \node[box=32mm,left=12mm of ckpt] (eval) {Đánh giá\\\scriptsize metrics, lỗi};

  \node[gbox=34mm,below=14mm of ckpt] (svc) {Service HTTP\\\scriptsize \code{/predict}, \code{/predict/batch}};
  \node[abox=34mm,below left=14mm and 2mm of svc] (ui1) {Giao diện: một câu};
  \node[abox=34mm,below right=14mm and 2mm of svc] (ui2) {Giao diện: dashboard};

  \draw[arr] (raw) -- (prep);
  \draw[arr] (prep) -- (proc);
  \draw[arr] (proc) -- (train);
  \draw[arr] (train) -- (ckpt);
  \draw[arr] (ckpt) -- (eval);
  \draw[arr] (ckpt) -- (svc) node[lab,midway,right] {nạp một lần};
  \draw[arr] (svc) -- (ui1);
  \draw[arr] (svc) -- (ui2);
\end{tikzpicture}
\caption{Ba tầng của hệ thống. Tầng nghiên cứu và tầng sản phẩm chỉ gặp nhau tại thư mục checkpoint và tại hợp đồng HTTP.}
\label{fig:kientruc}
\end{figure}

Lý do tách như vậy: tầng nghiên cứu phụ thuộc PyTorch và dữ liệu nặng, trong khi giao diện cần chạy được trên máy không
có GPU. Mã nạp mô hình được nhập khẩu \emph{bên trong hàm}, nên toàn bộ phần xử lý dữ liệu và đánh giá vẫn chạy được khi
máy không cài PyTorch --- đây cũng là điều cho phép bộ kiểm thử chạy nhanh trong môi trường phát triển.

\section{Ba giao thức chia tập}
\label{sec:giaothuc}

\begin{table}[H]
\centering
\caption{So sánh ba giao thức chia tập mà pipeline hỗ trợ.}
\label{tab:giaothuc}
\begin{tabularx}{\linewidth}{@{}lXl@{}}
\toprule
\textbf{Giao thức} & \textbf{Cách làm và tính chất} & \textbf{Kiểm tra rò rỉ} \\
\midrule
\code{clean\_split} & Đọc XML, chia \emph{theo bài đánh giá} với seed cố định & Có, bằng mã \\
\code{original\_split} & Dùng tệp chia sẵn, căn chỉnh từng dòng với câu trong XML để lấy lại định danh bài & Có, nếu căn chỉnh thành công \\
\code{provided\_split} & Đọc thẳng tệp \code{train/dev/test.txt}, mỗi dòng là một mẫu & Không thể: TXT không có định danh bài \\
\bottomrule
\end{tabularx}
\end{table}

Lượt chạy seed 42 dùng \code{provided\_split}. Lựa chọn này giữ đúng phân chia của tác giả bộ dữ liệu, nhưng đánh đổi:
bảo đảm \enquote{không rò rỉ ở mức bài đánh giá} không còn kiểm chứng được bằng mã. Hệ thống không che giấu đánh đổi đó
--- metadata ghi \code{review\_level\_leakage\_check.status = not\_verifiable}, và phần audit từ chối phát biểu
\code{no\_review\_leakage} cho dữ liệu dạng dòng.

\section{Hợp đồng HTTP}

Giao diện và mô hình chỉ gặp nhau qua hai điểm cuối. Dạng dữ liệu được chọn trùng với dạng mà \code{capstone.predict} đã
sinh ra, nên tầng phục vụ chỉ là lớp bọc mỏng chứ không phải một cách biểu diễn thứ hai.

\begin{Verbatim}[fontsize=\small,frame=single,framesep=3mm]
POST /predict
{"text": "The room was dirty.", "domain": "hotels", "model_type": "paraphrase"}

200 OK
{"text": ..., "domain": ..., "model_type": ...,
 "quadruples": [{"target": "room", "aspect": "rooms cleanliness",
                 "opinion": "dirty", "sentiment": "negative"}],
 "raw_generation": "...", "invalid_predictions": []}

POST /predict/batch
{"domain": "hotels", "model_type": "paraphrase",
 "items": [{"id": "row-2", "text": "..."}, {"id": "row-3", "text": "..."}]}

200 OK  {"results": [{"id": "row-2", "quadruples": [...]}, ...]}

Lỗi (mọi điểm cuối)
{"error": {"code": "MODEL_NOT_READY", "message": "..."}}
\end{Verbatim}

Ba quyết định đáng nói:

\begin{itemize}
  \item \textbf{Định danh do phía gọi đặt.} Mỗi phần tử trong lô mang \code{id} riêng và máy chủ trả lại đúng định danh
        đó. Phía giao diện kiểm tra tập định danh trả về khớp tập đã gửi, nên kết quả không thể gắn nhầm sang nhận xét
        khác khi thứ tự thay đổi.
  \item \textbf{Một dạng lỗi duy nhất.} Mọi lỗi, kể cả lỗi kiểm tra dữ liệu vào, đều mang dạng \code{\{error:\{code,
        message\}\}}; giao diện chỉ cần một nhánh xử lý.
  \item \textbf{Thiếu mô hình là một trạng thái, không phải sự cố.} Cặp (lĩnh vực, mô hình) chưa khai báo hoặc thiếu
        checkpoint trả về \code{MODEL\_NOT\_READY} kèm mã HTTP 503.
\end{itemize}

\section{Thiết kế dashboard theo khía cạnh}
\label{sec:thietkedashboard}

Phần này trả lời câu hỏi của doanh nghiệp: \emph{khía cạnh nào đang kéo đánh giá xuống?} Thiết kế xoay quanh một quyết
định duy nhất nhưng quan trọng: \textbf{đơn vị đếm là bài đánh giá, không phải lượt nhắc}.

Nếu đếm theo lượt nhắc, một bài chê phòng bẩn ba lần sẽ được tính ba lần và làm lệch tỷ lệ. Vì vậy mỗi bài đánh giá chỉ
đóng góp tối đa một lần cho mỗi khía cạnh, theo quy tắc gộp ở Bảng~\ref{tab:quytac}.

\begin{table}[H]
\centering
\caption{Quy tắc gộp nhiều lượt nhắc cùng một khía cạnh trong cùng một bài đánh giá.}
\label{tab:quytac}
\begin{tabularx}{\linewidth}{@{}lX@{}}
\toprule
\textbf{Nhãn bài} & \textbf{Điều kiện} \\
\midrule
Conflict & Có nhãn \code{conflict}, hoặc vừa có lượt tích cực vừa có lượt tiêu cực \\
Negative & Có ít nhất một lượt tiêu cực và không có lượt tích cực \\
Positive & Có ít nhất một lượt tích cực và không có lượt tiêu cực \\
Neutral  & Chỉ có các lượt trung tính \\
\bottomrule
\end{tabularx}
\end{table}

Quy tắc này khiến các tỷ lệ của một khía cạnh luôn cộng lại bằng 100\%, và giữ nhãn \code{conflict} tách riêng thay vì
gộp vào một nhóm khác --- phù hợp với việc OATS coi conflict là một nhãn độc lập. Quy tắc được in ngay trên giao diện
trong mục \enquote{How are reviews counted?} để người đọc bảng không phải đoán.

Ba lựa chọn thiết kế còn lại đều nhằm chống kết luận sai:

\begin{itemize}
  \item \textbf{Ngưỡng số mẫu tối thiểu} (mặc định 3 bài đánh giá) ẩn những khía cạnh quá ít dữ liệu, tránh hiển thị
        \enquote{100\% tiêu cực} dựa trên một nhận xét duy nhất.
  \item \textbf{Hai mức gom nhóm:} theo danh mục đầy đủ (\code{rooms cleanliness}) hoặc theo nhóm lớn (\code{rooms}),
        cho phép nhìn tổng quan trước rồi mới đi sâu.
  \item \textbf{Luôn kèm bằng chứng:} mỗi khía cạnh mở ra được danh sách cụm từ chê phổ biến và các nhận xét gốc, nên
        người dùng kiểm tra lại được thay vì tin vào một con số.
\end{itemize}

% ============================================================================
\chapter{Triển khai}
\label{chap:trienkhai}

\section{Tầng dữ liệu}

Bộ đọc XML kiểm tra tính nhất quán của chú thích ngay khi đọc: với mỗi quadruple, nếu dữ liệu cung cấp vị trí ký tự thì
vị trí đó phải trỏ đúng vào cụm từ tương ứng trong câu. Các bất thường không làm dừng chương trình mà được ghi thành
danh sách \code{DataIssue} có lý do, kèm đường dẫn nguồn và số dòng, rồi tổng hợp vào metadata.

Quyết định quan trọng ở bước làm sạch: \textbf{bài đánh giá không có nhãn được giữ lại}. Không có nhãn có thể nghĩa là
\enquote{câu này thật sự không chứa ý kiến} hoặc \enquote{người gán nhãn bỏ sót}; dữ liệu không phân biệt được hai
trường hợp đó. Vì vậy pipeline coi đây là trạng thái chưa rõ và chỉ cho phép loại bỏ thông qua một tệp kê khai có dẫn
chứng; nếu tệp kê khai chứa định danh không tồn tại, bước chuẩn bị dữ liệu dừng lại thay vì im lặng bỏ qua.

\section{Dấu vân tay và chống lệch dữ liệu}

Đây là cơ chế giữ cho kết quả không bị \enquote{trôi}. Bước chuẩn bị dữ liệu ghi vào \code{metadata.json}: mã băm
SHA-256 của từng tệp nguồn và từng tệp tập dữ liệu, cấu hình làm sạch, danh mục khía cạnh của lĩnh vực và, với
\code{clean\_split}, cả seed. Trước khi huấn luyện, trình huấn luyện đối chiếu metadata này với cấu hình thí nghiệm; khi
đánh giá, bước đánh giá kiểm tra lại dấu vân tay; khi so sánh hai mô hình, chương trình yêu cầu hai thí nghiệm có
\emph{cùng dấu vân tay tập test và cùng danh mục khía cạnh}, chứ không chỉ cùng số lượng mẫu.

Trong lần rà soát trước khi gộp mã, nhóm phát hiện metadata của giao thức \code{provided\_split} lưu đường dẫn tuyệt đối
của máy chạy. Hệ quả là dữ liệu đã xử lý không dùng lại được trên máy khác, vì bước kiểm tra sẽ không tìm thấy tệp
nguồn. Lỗi đã được sửa: đường dẫn nay được lưu đúng như cấu hình, thường là đường dẫn tương đối tới gốc dự án.

\section{Hai bộ chuyển đổi đầu ra}

Mỗi mô hình có một cặp hàm: một hàm biến danh sách quadruple thành chuỗi đích để huấn luyện, và một hàm đọc ngược chuỗi
sinh ra thành quadruple. Hàm đọc ngược phải chịu được mọi thứ mô hình có thể sinh ra, kể cả chuỗi hỏng, nên nó không
ném lỗi mà trả về hai danh sách: các quadruple hợp lệ, và các mẩu không đọc được kèm lý do. Các mẩu hỏng này đi thẳng
vào kết quả đánh giá dưới dạng \code{invalid\_predictions}, nhờ đó \enquote{mô hình sinh sai cú pháp} là một đại lượng
đo được chứ không phải thông tin bị mất.

\section{Chuẩn hóa và tính điểm}

Trước khi so khớp, cả nhãn chuẩn lẫn dự đoán đều đi qua bước chuẩn hóa: gộp khoảng trắng thừa, nhận diện các cách viết
của giá trị ngầm hiểu, và đối chiếu khía cạnh với danh mục của lĩnh vực. Việc đối chiếu danh mục quan trọng vì mô hình
sinh tự do có thể tạo ra một khía cạnh không tồn tại; những trường hợp đó được đánh dấu thay vì âm thầm tính là sai.

\subsection{Chỉ số nới lỏng cho dashboard}
\label{sec:relaxed}

Chỉ số chính của bài toán đòi đúng cả bốn thành phần, trong đó có hai chuỗi tự do là đối tượng và cụm ý kiến. Nhưng
dashboard ở Mục~\ref{sec:thietkedashboard} chỉ tiêu thụ cặp \emph{(khía cạnh, cảm xúc)}. Một dự đoán
\code{(bathroom, rooms cleanliness, \enquote{filthy}, negative)} cho nhãn chuẩn
\code{(room, rooms cleanliness, \enquote{dirty}, negative)} bị chỉ số chính tính là \emph{sai hoàn toàn}, trong khi với
dashboard nó \emph{đúng}: vẫn là một bài đánh giá tiêu cực về độ sạch phòng.

Nếu chỉ có chỉ số chính thì không có cách nào biết dashboard đáng tin đến đâu. Vì vậy bước đánh giá xuất thêm hai chỉ số
nới lỏng, tính bằng đúng cơ chế đếm chống trùng lặp của chỉ số chính nhưng sau khi chiếu quadruple xuống một tập con
các trường: chỉ \code{aspect}, và cặp \code{aspect + sentiment}. Trường \code{definition} đi kèm trong tệp JSON ghi
thẳng rằng đây là chỉ số nới lỏng và không được trình bày như kết quả ASQP --- để con số dễ nhìn hơn không bị trích dẫn
nhầm chỗ.

Ngoài P/R/F1 khớp hoàn toàn, bước đánh giá còn xuất các chẩn đoán: độ chính xác theo từng thành phần, kết quả theo số
lượng quadruple trong câu, và phân loại lỗi (thiếu quadruple, thừa quadruple, sai một thành phần, hỏng về đối tượng hoặc
ý kiến ngầm hiểu).

\begin{ghichu}{Một lỗi đếm trùng đã được sửa}
Khi rà soát, nhóm phát hiện một quadruple đã được báo cáo trong danh sách \enquote{sai một thành phần} vẫn bị đếm thêm
một lần nữa vào danh sách \enquote{hỏng ý kiến ngầm hiểu}. Nhánh xử lý đối tượng đã được khử trùng lặp trước đó, nhánh
ý kiến thì chưa. Lỗi này chỉ ảnh hưởng \emph{số liệu chẩn đoán}, không ảnh hưởng P/R/F1, và đã được sửa kèm ba test
kiểm chứng thất bại trên mã cũ. Con số \code{implicit\_opinion\_failure} công bố trong báo cáo lượt chạy seed 42 được
tính trước bản sửa này, nên không so sánh trực tiếp được với con số tương ứng của đối tượng.
\end{ghichu}

\section{Service dự đoán}

Service có hai phần. Phần đăng ký ánh xạ cặp (lĩnh vực, mô hình) tới một thư mục checkpoint theo tệp cấu hình, nạp mô
hình \emph{lười} ở lần gọi đầu tiên rồi giữ lại trong bộ nhớ. Phần ứng dụng chuyển hai điểm cuối thành lời gọi tới mô
hình đã nạp.

Quyết định đáng nói nhất là cách xử lý thất bại: mọi lỗi khi nạp mô hình --- thiếu thư viện, thiếu tệp, checkpoint hỏng
--- đều được ghi log đầy đủ rồi chuyển thành một \code{MODEL\_NOT\_READY} duy nhất gửi cho phía gọi, và kết quả thất bại
đó được nhớ lại để không thử nạp lại ở từng yêu cầu. Nhờ vậy một checkpoint hỏng không biến thành hàng nghìn lần nạp mô
hình thất bại.

\section{Giao diện web}

Giao diện tách thành ba lớp: các thành phần hiển thị, một giao diện lập trình \code{PredictionService}, và hai hiện thực
của giao diện đó --- một gọi HTTP thật, một chạy dữ liệu mẫu ngay trong trình duyệt. Thành phần hiển thị không biết hiện
thực nào đang chạy; chuyển đổi bằng biến môi trường, không phải bằng sửa mã.

Chế độ dữ liệu mẫu không phải là mã tạm thời. Nó cho phép phát triển và kiểm thử toàn bộ giao diện khi chưa có
checkpoint, và được dán nhãn rõ trên màn hình để không ai nhầm dữ liệu mẫu với dự đoán thật.

Hai chi tiết kỹ thuật đáng chú ý:

\begin{itemize}
  \item \textbf{Chống ghi đè kết quả cũ.} Mỗi yêu cầu mới hủy yêu cầu đang chạy; phản hồi đến muộn của yêu cầu đã bị
        thay thế sẽ bị bỏ qua, nên kết quả hiển thị luôn ứng với lần bấm cuối cùng.
  \item \textbf{Xử lý tệp theo lô, có thể tiếp tục.} Tệp nhiều nhận xét được chia thành các lô 20 dòng. Các lô đã xong
        được giữ lại, nên khi gặp lỗi mạng giữa chừng, nút \enquote{Resume} chỉ gửi lại phần còn thiếu thay vì làm lại
        từ đầu.
\end{itemize}

% ============================================================================
\chapter{Kết quả và đánh giá}
\label{chap:ketqua}

\section{Thiết lập thí nghiệm}

Sáu thí nghiệm (2 mô hình $\times$ 3 lĩnh vực) dùng chung thiết lập: T5-base, 20 epoch, kích thước lô 16, tốc độ học
$3\cdot10^{-4}$, độ dài chuỗi vào/ra 128, giải mã tham lam, seed 42, giao thức \code{provided\_split}. Checkpoint được
chọn theo F1 trên tập dev. Mọi số liệu dưới đây là của \emph{một} seed duy nhất.

\section{Kết quả trên tập test}

\begin{table}[H]
\centering
\caption{Kết quả seed 42 trên tập test, khớp hoàn toàn cả bốn thành phần.}
\label{tab:ketqua}
\begin{tabular}{llrrrrrrr}
\toprule
\textbf{Lĩnh vực} & \textbf{Mô hình} & \textbf{Epoch} & \textbf{Dev F1} & \textbf{P} & \textbf{R} & \textbf{F1} & \textbf{TP} & \textbf{FP} \\
\midrule
\code{amazon\_ff} & GAS        & 14 & 0,297 & 0,247 & 0,209 & \textbf{0,227} & 202 & 615 \\
\code{coursera}   & GAS        & 5  & 0,295 & 0,071 & 0,254 & \textbf{0,111} & 60  & 785 \\
\code{hotels}     & GAS        & 19 & 0,442 & 0,299 & 0,413 & \textbf{0,347} & 100 & 234 \\
\code{amazon\_ff} & Paraphrase & 5  & 0,294 & 0,215 & 0,205 & \textbf{0,210} & 198 & 721 \\
\code{coursera}   & Paraphrase & 3  & 0,308 & 0,095 & 0,309 & \textbf{0,145} & 73  & 699 \\
\code{hotels}     & Paraphrase & 16 & 0,454 & 0,337 & 0,459 & \textbf{0,389} & 111 & 218 \\
\bottomrule
\end{tabular}
\end{table}

So sánh hai mô hình trong cùng lĩnh vực: Paraphrase cao hơn GAS ở \code{coursera} ($+0{,}034$) và \code{hotels}
($+0{,}042$), thấp hơn ở \code{amazon\_ff} ($-0{,}017$). Ba quan sát từ một seed duy nhất không đủ để kết luận mô hình
nào tốt hơn nói chung; muốn phát biểu như vậy cần giao thức năm seed như bài báo.

\begin{figure}[H]
\centering
\begin{tikzpicture}
\begin{axis}[
  width=0.9\linewidth, height=6.2cm,
  ybar, bar width=9pt,
  ylabel={F1 khớp hoàn toàn},
  symbolic x coords={amazon\_ff, coursera, hotels},
  xtick=data, ymin=0, ymax=0.45,
  legend style={at={(0.5,-0.18)},anchor=north,legend columns=4,font=\small},
  ymajorgrids, grid style={lightgray},
  nodes near coords, nodes near coords style={font=\tiny},
  every node near coord/.append style={/pgf/number format/.cd,fixed,precision=3},
]
\addplot[fill=oatsblue!70,draw=oatsblue] coordinates {(amazon\_ff,0.227) (coursera,0.111) (hotels,0.347)};
\addplot[fill=oatsgreen!70,draw=oatsgreen] coordinates {(amazon\_ff,0.210) (coursera,0.145) (hotels,0.389)};
\addplot[fill=oatsgray!40,draw=oatsgray] coordinates {(amazon\_ff,0.196) (coursera,0.222) (hotels,0.263)};
\addplot[fill=oatsamber!50,draw=oatsamber] coordinates {(amazon\_ff,0.208) (coursera,0.198) (hotels,0.345)};
\legend{GAS (nhóm), Paraphrase (nhóm), GAS-T5 (bài báo), Paraphrase (bài báo)}
\end{axis}
\end{tikzpicture}
\caption{Kết quả của nhóm so với số công bố trong bài báo OATS. Hai nhóm cột \emph{không} cùng điều kiện thí nghiệm:
số của nhóm là một seed theo \code{provided\_split}, số của bài báo là trung bình năm seed; cú pháp đích của GAS cũng
khác. Hình chỉ để định vị độ lớn, không phải một phép so sánh công bằng.}
\label{fig:sosanh}
\end{figure}

\section{Vì sao không so sánh trực tiếp được với bài báo}
\label{sec:nhanrong}

Khi rà soát kết quả, nhóm nhận thấy \code{coursera} có 785 dự đoán sai dương so với chỉ 60 dự đoán đúng --- một tỷ lệ
bất thường. Nguyên nhân nằm ở dữ liệu chứ không ở mô hình. Bảng~\ref{tab:rong} cho thấy tỷ lệ dòng không mang nhãn nào
trong các tệp chia sẵn.

\begin{table}[H]
\centering
\caption{Tỷ lệ dòng không có nhãn trong các tệp chia sẵn, do nhóm đếm trực tiếp.}
\label{tab:rong}
\begin{tabular}{lrrr}
\toprule
\textbf{Lĩnh vực} & \textbf{Train} & \textbf{Dev} & \textbf{Test} \\
\midrule
\code{amazon\_ff} & 29,3\% & 23,5\% & 25,8\% \\
\code{coursera}   & 26,2\% & 23,9\% & \textbf{80,8\%} \\
\code{hotels}     & 24,5\% & 39,4\% & \textbf{86,7\%} \\
\bottomrule
\end{tabular}
\end{table}

Tổng số dòng và tổng số dòng không nhãn của ba tập cộng lại khớp với dữ liệu XML, nên đây đúng là phân chia gốc của bộ
dữ liệu chứ không phải lỗi xử lý. Nhưng hệ quả rất lớn: ở \code{hotels}, hơn 86\% câu trong tập test không có nhãn nào,
nên mỗi quadruple mà mô hình sinh ra cho những câu đó đều bị tính là sai dương, kéo precision xuống. Tỷ lệ này ở tập
train chỉ khoảng 25\%, tức là phân bố giữa lúc học và lúc chấm rất khác nhau.

\begin{ghichu}{Kết luận cần giữ đúng mức}
Các con số trong Bảng~\ref{tab:ketqua} là kết quả hợp lệ \emph{trong giao thức \code{provided\_split}}. Chúng không
tương đương với số của bài báo, và chênh lệch giữa các lĩnh vực phản ánh đặc điểm phân chia dữ liệu ít nhất ngang với
việc phản ánh chất lượng mô hình.
\end{ghichu}

\section{Trạng thái các giao thức chia tập}

Khi kiểm chứng lại, nhóm chạy thử cả ba giao thức trên dữ liệu thật:

\begin{table}[H]
\centering
\caption{Kết quả chạy thử bước chuẩn bị dữ liệu theo từng giao thức.}
\begin{tabular}{llll}
\toprule
\textbf{Giao thức} & \code{amazon\_ff} & \code{coursera} & \code{hotels} \\
\midrule
\code{clean\_split}    & Thành công & Thành công & Thành công \\
\code{provided\_split} & Thành công & Thành công & Thành công \\
\code{original\_split} & Lỗi (dòng 12) & Lỗi (dòng 11) & Lỗi (dòng 48) \\
\bottomrule
\end{tabular}
\end{table}

Lỗi của \code{original\_split} là \enquote{một dòng khớp với nhiều câu XML chưa được gán và nhãn không phân biệt được
duy nhất một câu}: khi cùng một câu xuất hiện nhiều lần trong bộ dữ liệu, bước căn chỉnh không biết dòng đó thuộc về
bài đánh giá nào. Phát hiện quan trọng là \code{clean\_split} --- giao thức có kiểm tra rò rỉ --- vẫn chạy được trên cả
ba lĩnh vực. Nói cách khác, giao thức chính của dự án không hỏng, và việc chạy lại ma trận thí nghiệm theo
\code{clean\_split} là khả thi. Repository giữ sẵn sáu cấu hình \code{clean\_split} với thư mục dữ liệu và thư mục kết
quả riêng để giao thức này không bị bỏ quên.

Dữ liệu cho giao thức này đã được chuẩn bị xong cho cả ba lĩnh vực, chia 80/10/10 ở mức bài đánh giá, và bước audit xác
nhận \code{no\_review\_leakage: true} ở cả ba. Phần còn thiếu chỉ là huấn luyện, vì lý do phần cứng ở
Mục~\ref{sec:phancung}.

\section{Hạ tầng huấn luyện và giới hạn phần cứng}
\label{sec:phancung}

Nhóm đã thử chạy ma trận \code{clean\_split} trên một máy xách tay có GPU RTX 5050 8\,GB VRAM và 16\,GB RAM. Lần thử
không hoàn tất, và các số đo thu được đáng ghi lại vì chúng quyết định việc chọn máy cho lần chạy sau.

\begin{table}[H]
\centering
\caption{Thời gian một bước huấn luyện, đo trên RTX 5050 8\,GB, batch 16, chuỗi vào/ra 128.}
\label{tab:tocdo}
\begin{tabular}{llr}
\toprule
\textbf{Mô hình} & \textbf{Precision} & \textbf{Giây/bước} \\
\midrule
Paraphrase & fp32 & 0,5 \\
GAS & fp32 & 3,8 \\
GAS & bf16 & 29,0 \\
\bottomrule
\end{tabular}
\end{table}

Ba điều rút ra:

\begin{itemize}
  \item \textbf{Nút thắt là VRAM, không phải độ dài chuỗi.} GAS chậm gấp 8 lần Paraphrase trên cùng mô hình và cùng
        kích thước lô. Mức chiếm VRAM đo được là 7{,}4--7{,}6\,GB trên card 8\,GB; phần vượt bị hệ điều hành đẩy sang
        RAM hệ thống, và tốc độ sụp theo.
  \item \textbf{Mixed precision không phải lúc nào cũng nhanh hơn.} Giả thiết ban đầu của nhóm là bật bf16 sẽ giảm bộ
        nhớ và tăng tốc. Đo thực tế cho kết quả ngược lại: bf16 chậm hơn fp32 tám lần, vì autocast vẫn giữ trọng số
        gốc ở fp32 rồi tạo thêm bản bf16 cho từng phép toán, tức \emph{tăng} đỉnh bộ nhớ trên một card vốn đã đầy.
        Tuỳ chọn \code{training.precision} vẫn được giữ lại trong mã, mặc định \code{fp32}, kèm cảnh báo này.
  \item \textbf{Hai lượt huấn luyện không được chạy song song.} Một lần chạy chồng nhau làm lượt đầu dừng với
        \code{CUDA error: CUBLAS\_STATUS\_EXECUTION\_FAILED} ở epoch 12 sau hơn 45 phút. Pipeline không hỗ trợ chạy
        tiếp dở dang, nên toàn bộ lượt đó phải làm lại.
\end{itemize}

Các bài học này được ghi thành hướng dẫn vận hành trong repository để lần chạy sau trên máy mạnh hơn không lặp lại,
trong đó có một bước bắt buộc: đo giây/bước trong vài phút đầu và dừng lại hỏi nếu tốc độ cho thấy tổng thời gian vượt
dự tính.

\section{Kiểm thử phần mềm}

\begin{table}[H]
\centering
\caption{Bộ kiểm thử tự động.}
\begin{tabularx}{\linewidth}{@{}lrX@{}}
\toprule
\textbf{Phạm vi} & \textbf{Số test} & \textbf{Nội dung chính} \\
\midrule
Pipeline (pytest)   & 133 & Đọc XML, làm sạch, chia tập và chống rò rỉ, dấu vân tay, hai bộ chuyển đổi đầu ra, chuẩn hóa và tính điểm (gồm chỉ số nới lỏng), đường dẫn cấu hình lồng thư mục, lựa chọn precision, các lệnh dòng lệnh \\
Service (pytest)    & 27  & Hợp đồng hai điểm cuối, echo định danh theo lô, nạp mô hình một lần, các đường lỗi \code{MODEL\_NOT\_READY} và \code{PREDICTION\_FAILED}, CORS, đọc tệp cấu hình phục vụ \\
Giao diện (Vitest)  & 98  & Trạng thái tải/lỗi/rỗng, hiển thị quadruple và giá trị ngầm hiểu, đọc tệp CSV/TXT/JSONL/JSON, quy tắc đếm theo bài đánh giá, xuất CSV, chia lô và tiếp tục khi lỗi \\
\bottomrule
\end{tabularx}
\end{table}

Ngoài test tự động, luồng web đã được chạy thật đầu-cuối: service chạy trên cổng 8000, giao diện chạy ở chế độ API và
gọi qua proxy, thao tác thật trên trình duyệt cho thấy nhãn chế độ đổi thành \enquote{API mode}, phân tích một câu trả
về quadruple đúng, cặp lĩnh vực chưa khai báo hiện \code{MODEL\_NOT\_READY} kèm nút thử lại, và tệp mẫu 30 nhận xét đi
qua điểm cuối theo lô rồi lên dashboard.

\begin{ghichu}{Điều kiện của lần kiểm chứng đầu-cuối}
Máy phát triển không cài PyTorch và không có checkpoint, nên lần chạy đầu-cuối dùng một bộ dự đoán giả đặt \emph{ngoài}
mã sản phẩm. Phép thử này xác nhận hợp đồng HTTP, việc chia lô và toàn bộ luồng giao diện; nó \emph{không} xác nhận chất
lượng dự đoán của mô hình thật.
\end{ghichu}

% ============================================================================
\chapter{Kết luận và hướng phát triển}
\label{chap:ketluan}

\section{Kết luận}

Hệ thống đã đi trọn một vòng: từ dữ liệu OATS thô, qua pipeline có kiểm chứng, tới sáu mô hình đã huấn luyện, và tới một
giao diện trả lời được câu hỏi kinh doanh \enquote{khía cạnh nào đang yếu và vì sao}. Phần nhóm cho là có giá trị nhất
không phải các con số F1 --- vốn còn thấp và chỉ từ một seed --- mà là việc mọi phát biểu trong hệ thống đều gắn với
bằng chứng: dấu vân tay dữ liệu chặn việc đánh giá nhầm tập, trạng thái không kiểm chứng được thì được ghi là không
kiểm chứng được, và chế độ dữ liệu mẫu được dán nhãn thay vì giả làm dự đoán thật.

\section{Hướng phát triển}

\begin{enumerate}
  \item \textbf{Lấy số cho cặp (khía cạnh, cảm xúc).} Chỉ số đã được cài đặt (Mục~\ref{sec:relaxed}) nhưng nó được tính
        lúc đánh giá, nên sáu lượt chạy cũ chưa có số này. Cần chạy lại bước đánh giá trên \code{predictions.jsonl} để
        biết dashboard đáng tin đến đâu trước khi đưa cho người dùng thật.
  \item \textbf{Chạy lại ma trận theo \code{clean\_split}} để có kết quả kèm bảo đảm không rò rỉ cấp bài đánh giá, rồi
        báo cáo song song hai giao thức. Dữ liệu đã sẵn sàng và đã kiểm tra không rò rỉ; chỉ còn phần huấn luyện, cần
        máy có VRAM lớn hơn theo Mục~\ref{sec:phancung}.
  \item \textbf{Giao thức năm seed} (30 lần huấn luyện, 30 lần đánh giá) để báo cáo trung bình và độ lệch chuẩn.
  \item \textbf{Sửa bước căn chỉnh XML} của \code{original\_split} cho trường hợp câu trùng lặp.
  \item \textbf{So sánh theo thời gian trên dashboard:} thêm cột ngày để doanh nghiệp kiểm chứng xem sau khi cải thiện
        một khía cạnh thì tỷ lệ tiêu cực có giảm thật hay không.
  \item \textbf{Tăng cường phần phục vụ:} xác thực, giới hạn tần suất và đóng gói container trước khi mở ra ngoài.
\end{enumerate}

\bibliographystyle{plain}
\bibliography{references}
\addcontentsline{toc}{chapter}{Tài liệu tham khảo}

\end{document}
